\subsubsection{Eventual Laguerre positivity on compact Fourier intervals}
\label{subsec:theta-high-order-compact}

A bounded-frequency interval imposes a substantially
weaker requirement than the spatial comparison on the
whole real line. In fact, the large-order part of the
generalized Laguerre hierarchy can be controlled without
assuming RH or any inequality involving an unknown
coefficient $c_n(s,x)$.

\begin{theorem}[Uniform high-order positivity on compact intervals]
\label{thm:compact-eventual-laguerre}
Let $H_t$ be the genuine Riemann theta deformation
\eqref{eq:heatH}, with $0\le t\le T=1/2$. For every
$X<\infty$ there exists an integer $N_X$ such that
\begin{equation}
\boxed{L_{n,t}(x)>0
\quad(n\ge N_X,\ |x|\le X,\ 0\le t\le T).}
\label{eq:compact-eventual}
\end{equation}
More precisely, the normalized Fourier pair measure
defined below satisfies
\begin{equation}
\boxed{\lim_{n\to\infty}\ \sup_{0\le t\le T}
\mathbb E_{n,t}(u+v)^2=0.}
\label{eq:conditional-squeeze}
\end{equation}
This is a compact-frequency, large-order assertion;
the index $N_X$ is not claimed to be independent of $X$.
\end{theorem}

\begin{proof}
Put $h_t(u)=\tfrac12e^{tu^2}\Phi(u)$, an even,
strictly positive, rapidly decreasing function, and
$V(u)=-\log\Phi(u)$. The Fourier product identity yields
\begin{equation}
L_{n,t}(x)=\frac{1}{(2n)!}
\iint_{\mathbb R^2}
h_t(u)h_t(v)(u-v)^{2n}
\cos(x(u+v))\,du\,dv .
\label{eq:high-order-pair}
\end{equation}
In particular $A_{n,t}=L_{n,t}(0)>0$. Let
$\mathbb P_{n,t}$ be the probability measure obtained
by dividing the integrand at $x=0$ by its total mass,
and write $W=u+v$, $D=u-v$. We have
\begin{equation}
\frac{L_{n,t}(x)}{A_{n,t}}
=\mathbb E_{n,t}\cos(xW)
\ge1-\frac{x^2}{2}\mathbb E_{n,t}W^2 .
\label{eq:compact-cos-gate}
\end{equation}

We estimate the last expectation. Conditional on
$D=d$, the variable $W$ has a density proportional to
$e^{-Q_{d,t}(w)}$, where
\[
Q_{d,t}(w)
=V\!\left(\frac{w+d}{2}\right)
+V\!\left(\frac{w-d}{2}\right)
-\frac t2(w^2+d^2).
\]
The density is even in $w$, and
\begin{equation}
Q_{d,t}''(w)
=\frac14\left[
V''\!\left(\frac{w+d}{2}\right)
+V''\!\left(\frac{w-d}{2}\right)-4t
\right].
\label{eq:conditional-curvature}
\end{equation}
Newman's theta-potential curvature theorem
\cite{NewmanGHS} gives that $V''$ is even and
nondecreasing in $|u|$. Its first-summand
asymptotics also give
$V''(u)\sim16\pi e^{4u}$ as $u\to+\infty$.
Let $m=V''(0)$. Since at least one of
$|(w+d)/2|$ and $|(w-d)/2|$ is at least
$|d|/2$, for $|d|\ge D$ we obtain
\[
Q_{d,t}''(w)\ge
\kappa_D:=
\frac{V''(D/2)+m-2}{4}.
\]
In particular $\kappa_D\to\infty$.
For sufficiently large $D$ it is positive.
The elementary integration-by-parts identity
\[
1=\mathbb E(WQ_{d,t}'(W)\mid D=d)
\]
and $wQ_{d,t}'(w)\ge\kappa_Dw^2$
then show
\begin{equation}
\mathbb E(W^2\mid D=d)\le\kappa_D^{-1}
\quad(|d|\ge D).
\label{eq:conditional-variance}
\end{equation}

For the remaining compact set $|d|\le D$,
the conditional second moment is bounded,
uniformly in $t\in[0,T]$, by some finite
$B_D$. This follows directly from the
positive continuous denominator and the
uniform superexponential tails in the
conditional $w$-integral.
Moreover the factor $|u-v|^{2n}$ drives
the separation away from every bounded
set, uniformly in heat time.
Indeed put
$H_T=\int_{\mathbb R}h_T(u)\,du<\infty$
and choose a fixed rectangle
$\mathcal R_D$ on which $|u-v|\ge D+1$.
The number
$b_D=\iint_{\mathcal R_D}h_0(u)h_0(v)\,du\,dv$
is positive. Because $h_0\le h_t\le h_T$,
\begin{equation}
\mathbb P_{n,t}(|D|\le D)
\le \frac{H_T^2}{b_D}
\left(\frac{D}{D+1}\right)^{2n}.
\label{eq:separation-concentration}
\end{equation}
Equations \eqref{eq:conditional-variance} and
\eqref{eq:separation-concentration} imply
\begin{equation}
\sup_{0\le t\le T}\mathbb E_{n,t}W^2
\le B_D\frac{H_T^2}{b_D}
\left(\frac{D}{D+1}\right)^{2n}
+\kappa_D^{-1}.
\label{eq:highorder-second-moment}
\end{equation}
First let $n\to\infty$ and then
$D\to\infty$ to prove
\eqref{eq:conditional-squeeze}.
Finally choose $n$ sufficiently large
that the right side of
\eqref{eq:compact-cos-gate} is strictly
positive for all $|x|\le X$, uniformly
in $t$. This proves
\eqref{eq:compact-eventual}.
\end{proof}

\begin{corollary}[Finite-order reduction on a fixed frequency slab]
\label{cor:compact-finite-order}
For every finite $X$ the task of proving
\[
L_{n,t}(x)\ge0\quad
(n\ge1,\ 0\le t\le1/2,\ |x|\le X)
\]
reduces to the finitely many orders
$1\le n<N_X$. Moreover, for each
fixed $n\ge N_X$, the strictly positive
denominator $U_n(s,x)$ permits a
nonnegative, locally bounded spatial
comparison coefficient on
$[0,1/2]\times[-X,X]$.
\end{corollary}
\begin{proof}
The first assertion is
Theorem~\ref{thm:compact-eventual-laguerre}.
For the second, $U_n$ and $U_{n+1}$
are continuous and strictly positive on
the stated compact set, so their quotient
is bounded there. One may take
$c_n=k_nU_{n+1}/U_n$.
This construction does not assume the
comparison to prove the denominator's
positivity: that positivity was already
established independently above.
\end{proof}

\paragraph{Scope and connection with prior literature.}
The proof is based on concentration of
a positive theta pair measure, not on a
finite mesh of Laguerre values. It is
analogous in spirit, but not identical
in its indexing or formulation, to
eventual Jensen-polynomial hyperbolicity
in the derivative aspect
\cite{GORZ,JensenXi}.
In particular, no novelty priority is claimed
without a more complete comparison of
these two formulations.
The theorem provides no explicit numerical
value of $N_X$; proving an effective
joint rate in $n$ and $X$ is the next
quantitative problem. A fixed finite $X$
cannot be replaced by all of $\mathbb R$
without an additional argument.
The analytic integration-by-parts and
separation-concentration steps have not
been formalized end-to-end in Lean.
The associated finite algebraic gates are
recorded separately in the repository.
